\chapter{PHÂN TÍCH HÀNH VI}

\section{Phân tích hành vi bằng biểu đồ trình tự}

Sau phân tích chức năng (use case) và cấu trúc (lớp), chương này dùng \textbf{biểu đồ trình tự (Sequence)} để mô tả các đối tượng tương tác theo thời gian khi thực hiện use case.

\textbf{Nguyên tắc}: \textbf{mỗi use case đã đặc tả ở Ch.3 (UC01--UC34) có đúng một biểu đồ trình tự}. Các sơ đồ dùng chung mẫu lifeline: \textit{Actor} -- \textit{Giao diện (boundary)} -- \textit{Hệ thống (control)} -- \textit{CSDL (entity)}. CRUD cùng entity khác nhau chủ yếu ở tên thông điệp/phương thức và ràng buộc (xóa vs ngừng, kiểm tra tồn, trạng thái phiếu\ldots).

%---------- UC01 ----------
\subsection{UC01 -- Đăng nhập}

\textbf{Bối cảnh}: Xác thực để vào hệ thống. \textbf{Đối tượng}: Người dùng, giao diện đăng nhập, hệ thống, CSDL (\texttt{NguoiDung}).

\input{diagrams/tex/5.7_sequence_uc01}

Giải thích: Actor nhập form $\rightarrow$ hệ thống kiểm tra CSDL $\rightarrow$ tạo phiên hoặc báo lỗi. Luồng boundary $\rightarrow$ control $\rightarrow$ entity.

%---------- UC02 ----------
\subsection{UC02 -- Đăng xuất}

\textbf{Bối cảnh}: Kết thúc phiên làm việc. \textbf{Đối tượng}: Người dùng, giao diện, hệ thống, CSDL (phiên / log).

\begin{figure}[H]
\centering
\resizebox{\textwidth}{!}{%
\begin{sequencediagram}
\newinst{a}{Người dùng}
\newinst[2]{gd}{Giao diện hệ thống}
\newinst[2]{ht}{Hệ thống}
\newinst[2]{db}{Cơ sở dữ liệu}

\begin{call}{a}{1. Chọn đăng xuất}{gd}{}
\end{call}

\begin{call}{gd}{2. Gửi yêu cầu đăng xuất}{ht}{}
  \begin{call}{ht}{3. Hủy phiên làm việc}{db}{Xác nhận}
  \end{call}
  \begin{call}{ht}{4. Ghi log (tùy chọn)}{db}{}
  \end{call}
\end{call}

\begin{call}{ht}{5. Trả kết quả}{gd}{Đã đăng xuất}
\end{call}

\begin{call}{gd}{6. Chuyển trang đăng nhập}{a}{}
\end{call}

\end{sequencediagram}
}
\caption{Sequence diagram UC02 -- Đăng xuất}
\label{fig:sequence-uc02}
\end{figure}


Giải thích: Hủy phiên (không kiểm tra mật khẩu); có thể ghi log; chuyển trang đăng nhập.

%---------- UC03 ----------
\subsection{UC03 -- Thêm tài khoản}

\textbf{Bối cảnh}: QTV tạo tài khoản mới, gán vai trò. \textbf{Đối tượng}: QTV, giao diện tài khoản, hệ thống, CSDL.

\begin{figure}[H]
\centering
\resizebox{\textwidth}{!}{%
\begin{sequencediagram}
\newinst{qtv}{Quản trị viên}
\newinst[2]{gd}{Giao diện tài khoản}
\newinst[2]{ht}{Hệ thống}
\newinst[2]{db}{Cơ sở dữ liệu}

\begin{call}{qtv}{1. Chọn thêm tài khoản}{gd}{Hiển thị form}
\end{call}

\begin{call}{qtv}{2. Nhập thông tin, vai trò}{gd}{}
\end{call}

\begin{call}{gd}{3. Gửi thông tin tài khoản}{ht}{}
  \begin{call}{ht}{4. Kiểm tra trùng tên đăng nhập}{db}{Kết quả}
  \end{call}
  \begin{call}{ht}{5. Mã hóa mật khẩu, lưu NguoiDung}{db}{Xác nhận}
  \end{call}
  \begin{call}{ht}{6. Ghi log (tùy chọn)}{db}{}
  \end{call}
\end{call}

\begin{call}{ht}{7. Trả kết quả}{gd}{Thông báo thành công}
\end{call}

\begin{call}{gd}{8. Làm mới danh sách tài khoản}{qtv}{}
\end{call}

\end{sequencediagram}
}
\caption{Sequence diagram UC03 -- Thêm tài khoản}
\label{fig:sequence-uc03}
\end{figure}


Giải thích: Kiểm tra trùng tên đăng nhập $\rightarrow$ mã hóa mật khẩu, lưu \texttt{NguoiDung} $\rightarrow$ làm mới danh sách.

%---------- UC04 ----------
\subsection{UC04 -- Sửa tài khoản}

\textbf{Bối cảnh}: QTV cập nhật họ tên, email, vai trò, (mật khẩu nếu đổi).

\begin{figure}[H]
\centering
\resizebox{\textwidth}{!}{%
\begin{sequencediagram}
\newinst{a}{Quản trị viên}
\newinst[2]{gd}{Giao diện tài khoản}
\newinst[2]{ht}{Hệ thống}
\newinst[2]{db}{Cơ sở dữ liệu}

\begin{call}{a}{1. Chọn tài khoản cần sửa}{gd}{Hiển thị form}
\end{call}

\begin{call}{a}{2. Chỉnh sửa thông tin / vai trò}{gd}{}
\end{call}

\begin{call}{gd}{3. Gửi thông tin cập nhật}{ht}{}
  \begin{call}{ht}{4. Kiểm tra dữ liệu hợp lệ}{db}{Kết quả}
  \end{call}
  \begin{call}{ht}{5. Cập nhật NguoiDung}{db}{Xác nhận}
  \end{call}
\end{call}

\begin{call}{ht}{6. Trả kết quả}{gd}{Thông báo thành công}
\end{call}

\begin{call}{gd}{7. Làm mới danh sách tài khoản}{a}{}
\end{call}

\end{sequencediagram}
}
\caption{Sequence diagram UC04 -- Sửa tài khoản}
\label{fig:sequence-uc04}
\end{figure}


Giải thích: Chọn bản ghi $\rightarrow$ form sửa $\rightarrow$ validate $\rightarrow$ cập nhật \texttt{NguoiDung}.

%---------- UC05 ----------
\subsection{UC05 -- Xóa tài khoản}

\textbf{Bối cảnh}: QTV xóa hoặc khóa tài khoản tùy ràng buộc lịch sử/phiếu.

\begin{figure}[H]
\centering
\resizebox{\textwidth}{!}{%
\begin{sequencediagram}
\newinst{a}{Quản trị viên}
\newinst[2]{gd}{Giao diện tài khoản}
\newinst[2]{ht}{Hệ thống}
\newinst[2]{db}{Cơ sở dữ liệu}

\begin{call}{a}{1. Chọn xóa / khóa tài khoản}{gd}{Xác nhận}
\end{call}

\begin{call}{gd}{2. Gửi yêu cầu xóa/khóa}{ht}{}
  \begin{call}{ht}{3. Kiểm tra ràng buộc (phiếu, log)}{db}{Kết quả}
  \end{call}
  \begin{call}{ht}{4. Xóa hoặc khóa NguoiDung}{db}{Xác nhận}
  \end{call}
\end{call}

\begin{call}{ht}{5. Trả kết quả}{gd}{Thông báo}
\end{call}

\begin{call}{gd}{6. Cập nhật danh sách tài khoản}{a}{}
\end{call}

\end{sequencediagram}
}
\caption{Sequence diagram UC05 -- Xóa tài khoản}
\label{fig:sequence-uc05}
\end{figure}


Giải thích: Kiểm tra ràng buộc; nếu đã phát sinh dữ liệu thì \textbf{khóa} thay vì xóa cứng.

%---------- UC06 ----------
\subsection{UC06 -- Tìm kiếm tài khoản}

\textbf{Bối cảnh}: QTV lọc theo tên đăng nhập, họ tên, vai trò, trạng thái.

\begin{figure}[H]
\centering
\resizebox{\textwidth}{!}{%
\begin{sequencediagram}
\newinst{a}{Quản trị viên}
\newinst[2]{gd}{Giao diện tài khoản}
\newinst[2]{ht}{Hệ thống}
\newinst[2]{db}{Cơ sở dữ liệu}

\begin{call}{a}{1. Nhập từ khóa / bộ lọc}{gd}{}
\end{call}

\begin{call}{gd}{2. Gửi điều kiện tìm kiếm}{ht}{}
  \begin{call}{ht}{3. Truy vấn NguoiDung}{db}{Danh sách}
  \end{call}
\end{call}

\begin{call}{ht}{4. Trả kết quả}{gd}{Danh sách khớp}
\end{call}

\begin{call}{gd}{5. Hiển thị kết quả tìm kiếm}{a}{}
\end{call}

\end{sequencediagram}
}
\caption{Sequence diagram UC06 -- Tìm kiếm tài khoản}
\label{fig:sequence-uc06}
\end{figure}


Giải thích: Luồng \textbf{chỉ đọc} -- gửi điều kiện, truy vấn, hiển thị kết quả.

%---------- UC07 ----------
\subsection{UC07 -- Xem danh sách tài khoản}

\textbf{Bối cảnh}: QTV mở màn hình quản lý, tải danh sách (phân trang).

\begin{figure}[H]
\centering
\resizebox{\textwidth}{!}{%
\begin{sequencediagram}
\newinst{a}{Quản trị viên}
\newinst[2]{gd}{Giao diện tài khoản}
\newinst[2]{ht}{Hệ thống}
\newinst[2]{db}{Cơ sở dữ liệu}

\begin{call}{a}{1. Mở quản lý tài khoản}{gd}{}
\end{call}

\begin{call}{gd}{2. Yêu cầu danh sách}{ht}{}
  \begin{call}{ht}{3. Lấy danh sách NguoiDung}{db}{Dữ liệu}
  \end{call}
\end{call}

\begin{call}{ht}{4. Trả danh sách}{gd}{}
\end{call}

\begin{call}{gd}{5. Hiển thị bảng tài khoản}{a}{}
\end{call}

\end{sequencediagram}
}
\caption{Sequence diagram UC07 -- Xem danh sách tài khoản}
\label{fig:sequence-uc07}
\end{figure}


Giải thích: Tải toàn bộ/phân trang \texttt{NguoiDung}; điểm vào cho UC03--UC06.

%---------- UC08 ----------
\subsection{UC08 -- Thêm nhà cung cấp}

\textbf{Bối cảnh}: NV/QL bổ sung NCC phục vụ phiếu nhập.

\begin{figure}[H]
\centering
\resizebox{\textwidth}{!}{%
\begin{sequencediagram}
\newinst{nv}{Nhân viên / QL kho}
\newinst[2]{gd}{Giao diện NCC}
\newinst[2]{ht}{Hệ thống}
\newinst[2]{db}{Cơ sở dữ liệu}

\begin{call}{nv}{1. Chọn thêm nhà cung cấp}{gd}{Hiển thị form}
\end{call}

\begin{call}{nv}{2. Nhập thông tin NCC}{gd}{}
\end{call}

\begin{call}{gd}{3. Gửi thông tin NCC}{ht}{}
  \begin{call}{ht}{4. Kiểm tra dữ liệu / trùng mã}{db}{Kết quả}
  \end{call}
  \begin{call}{ht}{5. Lưu NhaCungCap}{db}{Xác nhận}
  \end{call}
\end{call}

\begin{call}{ht}{6. Trả kết quả}{gd}{Thông báo thành công}
\end{call}

\begin{call}{gd}{7. Hiển thị danh sách NCC}{nv}{}
\end{call}

\end{sequencediagram}
}
\caption{Sequence diagram UC08 -- Thêm nhà cung cấp}
\label{fig:sequence-uc08}
\end{figure}


Giải thích: Validate $\rightarrow$ insert \texttt{NhaCungCap}; không đụng tồn kho.

%---------- UC09 ----------
\subsection{UC09 -- Sửa nhà cung cấp}

\textbf{Bối cảnh}: Cập nhật thông tin liên hệ, địa chỉ, trạng thái NCC.

\begin{figure}[H]
\centering
\resizebox{\textwidth}{!}{%
\begin{sequencediagram}
\newinst{a}{Nhân viên / QL kho}
\newinst[2]{gd}{Giao diện NCC}
\newinst[2]{ht}{Hệ thống}
\newinst[2]{db}{Cơ sở dữ liệu}

\begin{call}{a}{1. Chọn NCC cần sửa}{gd}{Hiển thị form}
\end{call}

\begin{call}{a}{2. Chỉnh sửa thông tin NCC}{gd}{}
\end{call}

\begin{call}{gd}{3. Gửi thông tin cập nhật}{ht}{}
  \begin{call}{ht}{4. Kiểm tra dữ liệu}{db}{Kết quả}
  \end{call}
  \begin{call}{ht}{5. Cập nhật NhaCungCap}{db}{Xác nhận}
  \end{call}
\end{call}

\begin{call}{ht}{6. Trả kết quả}{gd}{Thành công}
\end{call}

\begin{call}{gd}{7. Làm mới danh sách NCC}{a}{}
\end{call}

\end{sequencediagram}
}
\caption{Sequence diagram UC09 -- Sửa nhà cung cấp}
\label{fig:sequence-uc09}
\end{figure}


Giải thích: Chọn NCC $\rightarrow$ cập nhật thuộc tính $\rightarrow$ làm mới danh sách.

%---------- UC10 ----------
\subsection{UC10 -- Xóa nhà cung cấp}

\textbf{Bối cảnh}: Xóa NCC; nếu đã có phiếu nhập thì chỉ ngừng hoạt động.

\begin{figure}[H]
\centering
\resizebox{\textwidth}{!}{%
\begin{sequencediagram}
\newinst{a}{Nhân viên / QL kho}
\newinst[2]{gd}{Giao diện NCC}
\newinst[2]{ht}{Hệ thống}
\newinst[2]{db}{Cơ sở dữ liệu}

\begin{call}{a}{1. Chọn xóa NCC}{gd}{Xác nhận}
\end{call}

\begin{call}{gd}{2. Gửi yêu cầu xóa}{ht}{}
  \begin{call}{ht}{3. Kiểm tra đã có phiếu nhập?}{db}{Kết quả}
  \end{call}
  \begin{call}{ht}{4. Xóa hoặc ngừng hoạt động}{db}{Xác nhận}
  \end{call}
\end{call}

\begin{call}{ht}{5. Trả kết quả}{gd}{Thông báo}
\end{call}

\begin{call}{gd}{6. Cập nhật danh sách NCC}{a}{}
\end{call}

\end{sequencediagram}
}
\caption{Sequence diagram UC10 -- Xóa nhà cung cấp}
\label{fig:sequence-uc10}
\end{figure}


Giải thích: Kiểm tra ràng buộc phiếu nhập trước khi xóa cứng / ngừng.

%---------- UC11 ----------
\subsection{UC11 -- Tìm kiếm nhà cung cấp}

\textbf{Bối cảnh}: Tra cứu NCC theo mã, tên, SĐT.

\begin{figure}[H]
\centering
\resizebox{\textwidth}{!}{%
\begin{sequencediagram}
\newinst{a}{Nhân viên / QL kho}
\newinst[2]{gd}{Giao diện NCC}
\newinst[2]{ht}{Hệ thống}
\newinst[2]{db}{Cơ sở dữ liệu}

\begin{call}{a}{1. Nhập từ khóa tìm NCC}{gd}{}
\end{call}

\begin{call}{gd}{2. Gửi điều kiện tìm}{ht}{}
  \begin{call}{ht}{3. Truy vấn NhaCungCap}{db}{Danh sách}
  \end{call}
\end{call}

\begin{call}{ht}{4. Trả kết quả}{gd}{}
\end{call}

\begin{call}{gd}{5. Hiển thị NCC khớp}{a}{}
\end{call}

\end{sequencediagram}
}
\caption{Sequence diagram UC11 -- Tìm kiếm nhà cung cấp}
\label{fig:sequence-uc11}
\end{figure}


Giải thích: Query chỉ đọc trên \texttt{NhaCungCap}.

%---------- UC12 ----------
\subsection{UC12 -- Xem danh sách nhà cung cấp}

\textbf{Bối cảnh}: Mở danh mục NCC, xem bảng đầy đủ.

\begin{figure}[H]
\centering
\resizebox{\textwidth}{!}{%
\begin{sequencediagram}
\newinst{a}{Nhân viên / QL kho}
\newinst[2]{gd}{Giao diện NCC}
\newinst[2]{ht}{Hệ thống}
\newinst[2]{db}{Cơ sở dữ liệu}

\begin{call}{a}{1. Mở danh mục NCC}{gd}{}
\end{call}

\begin{call}{gd}{2. Yêu cầu danh sách}{ht}{}
  \begin{call}{ht}{3. Lấy danh sách NhaCungCap}{db}{Dữ liệu}
  \end{call}
\end{call}

\begin{call}{ht}{4. Trả danh sách}{gd}{}
\end{call}

\begin{call}{gd}{5. Hiển thị bảng NCC}{a}{}
\end{call}

\end{sequencediagram}
}
\caption{Sequence diagram UC12 -- Xem danh sách nhà cung cấp}
\label{fig:sequence-uc12}
\end{figure}


Giải thích: \texttt{layDanhSach} NCC; điểm vào CRUD NCC.

%---------- UC13 ----------
\subsection{UC13 -- Thêm hàng hóa}

\textbf{Bối cảnh}: Thêm mặt hàng; khởi tạo tồn = 0 tại các kho.

\input{diagrams/tex/5.1_sequence_uc05}

Giải thích: Sau validate: lưu \texttt{HangHoa} rồi tạo \texttt{TonKho} ban đầu -- hai thao tác CSDL tuần tự.

%---------- UC14 ----------
\subsection{UC14 -- Sửa hàng hóa}

\textbf{Bối cảnh}: Cập nhật tên, giá, nhóm, mô tả (không tạo tồn mới).

\begin{figure}[H]
\centering
\resizebox{\textwidth}{!}{%
\begin{sequencediagram}
\newinst{a}{Nhân viên / QL kho}
\newinst[2]{gd}{Giao diện hàng hóa}
\newinst[2]{ht}{Hệ thống}
\newinst[2]{db}{Cơ sở dữ liệu}

\begin{call}{a}{1. Chọn hàng cần sửa}{gd}{Hiển thị form}
\end{call}

\begin{call}{a}{2. Chỉnh sửa thông tin / giá}{gd}{}
\end{call}

\begin{call}{gd}{3. Gửi thông tin cập nhật}{ht}{}
  \begin{call}{ht}{4. Kiểm tra dữ liệu}{db}{Kết quả}
  \end{call}
  \begin{call}{ht}{5. Cập nhật HangHoa}{db}{Xác nhận}
  \end{call}
\end{call}

\begin{call}{ht}{6. Trả kết quả}{gd}{Thành công}
\end{call}

\begin{call}{gd}{7. Làm mới danh sách hàng}{a}{}
\end{call}

\end{sequencediagram}
}
\caption{Sequence diagram UC14 -- Sửa hàng hóa}
\label{fig:sequence-uc14}
\end{figure}


Giải thích: Cập nhật \texttt{HangHoa}; tồn hiện có giữ nguyên trừ khi nghiệp vụ khác điều chỉnh.

%---------- UC15 ----------
\subsection{UC15 -- Xóa hàng hóa}

\textbf{Bối cảnh}: Xóa hoặc ngừng dùng hàng nếu đã phát sinh phiếu.

\input{diagrams/tex/5.seq_uc15}

Giải thích: Kiểm tra phiếu/tồn liên quan trước khi xóa cứng hoặc ngừng.

%---------- UC16 ----------
\subsection{UC16 -- Tìm kiếm hàng hóa}

\textbf{Bối cảnh}: Tìm theo mã/tên/nhóm; có thể hiển thị kèm tồn.

\input{diagrams/tex/5.2_sequence_uc06}

Giải thích: Luồng chỉ đọc; có thể ghép \texttt{TonKho} khi trả kết quả.

%---------- UC17 ----------
\subsection{UC17 -- Xem danh sách hàng hóa}

\textbf{Bối cảnh}: Xem danh mục hàng (phân trang), điểm vào CRUD hàng.

\begin{figure}[H]
\centering
\resizebox{\textwidth}{!}{%
\begin{sequencediagram}
\newinst{a}{Nhân viên / QL kho}
\newinst[2]{gd}{Giao diện hàng hóa}
\newinst[2]{ht}{Hệ thống}
\newinst[2]{db}{Cơ sở dữ liệu}

\begin{call}{a}{1. Mở danh mục hàng hóa}{gd}{}
\end{call}

\begin{call}{gd}{2. Yêu cầu danh sách}{ht}{}
  \begin{call}{ht}{3. Lấy HangHoa (+ tồn tóm tắt)}{db}{Dữ liệu}
  \end{call}
\end{call}

\begin{call}{ht}{4. Trả danh sách}{gd}{}
\end{call}

\begin{call}{gd}{5. Hiển thị bảng hàng hóa}{a}{}
\end{call}

\end{sequencediagram}
}
\caption{Sequence diagram UC17 -- Xem danh sách hàng hóa}
\label{fig:sequence-uc17}
\end{figure}


Giải thích: Tải \texttt{HangHoa} (+ tóm tắt tồn nếu cần).

%---------- UC18 ----------
\subsection{UC18 -- Thêm kho}

\textbf{Bối cảnh}: QL kho tạo kho mới (QTV không thực hiện nhóm kho).

\begin{figure}[H]
\centering
\resizebox{\textwidth}{!}{%
\begin{sequencediagram}
\newinst{ql}{Quản lý kho}
\newinst[2]{gd}{Giao diện kho}
\newinst[2]{ht}{Hệ thống}
\newinst[2]{db}{Cơ sở dữ liệu}

\begin{call}{ql}{1. Chọn thêm kho}{gd}{Hiển thị form}
\end{call}

\begin{call}{ql}{2. Nhập mã, tên, địa chỉ, phụ trách}{gd}{}
\end{call}

\begin{call}{gd}{3. Gửi thông tin kho}{ht}{}
  \begin{call}{ht}{4. Kiểm tra trùng mã kho}{db}{Kết quả}
  \end{call}
  \begin{call}{ht}{5. Lưu Kho}{db}{Xác nhận}
  \end{call}
\end{call}

\begin{call}{ht}{6. Trả kết quả}{gd}{Thông báo thành công}
\end{call}

\begin{call}{gd}{7. Hiển thị danh sách kho}{ql}{}
\end{call}

\end{sequencediagram}
}
\caption{Sequence diagram UC18 -- Thêm kho}
\label{fig:sequence-uc18}
\end{figure}


Giải thích: Kiểm tra trùng mã $\rightarrow$ lưu \texttt{Kho}.

%---------- UC19 ----------
\subsection{UC19 -- Sửa kho}

\textbf{Bối cảnh}: Cập nhật tên, địa chỉ, người phụ trách, trạng thái kho.

\begin{figure}[H]
\centering
\resizebox{\textwidth}{!}{%
\begin{sequencediagram}
\newinst{a}{Quản lý kho}
\newinst[2]{gd}{Giao diện kho}
\newinst[2]{ht}{Hệ thống}
\newinst[2]{db}{Cơ sở dữ liệu}

\begin{call}{a}{1. Chọn kho cần sửa}{gd}{Hiển thị form}
\end{call}

\begin{call}{a}{2. Chỉnh sửa thông tin kho}{gd}{}
\end{call}

\begin{call}{gd}{3. Gửi thông tin cập nhật}{ht}{}
  \begin{call}{ht}{4. Kiểm tra dữ liệu}{db}{Kết quả}
  \end{call}
  \begin{call}{ht}{5. Cập nhật Kho}{db}{Xác nhận}
  \end{call}
\end{call}

\begin{call}{ht}{6. Trả kết quả}{gd}{Thành công}
\end{call}

\begin{call}{gd}{7. Làm mới danh sách kho}{a}{}
\end{call}

\end{sequencediagram}
}
\caption{Sequence diagram UC19 -- Sửa kho}
\label{fig:sequence-uc19}
\end{figure}


Giải thích: Cập nhật thuộc tính \texttt{Kho}.

%---------- UC20 ----------
\subsection{UC20 -- Xóa kho}

\textbf{Bối cảnh}: Xóa/ngừng kho; chặn nếu còn tồn hoặc phiếu.

\begin{figure}[H]
\centering
\resizebox{\textwidth}{!}{%
\begin{sequencediagram}
\newinst{a}{Quản lý kho}
\newinst[2]{gd}{Giao diện kho}
\newinst[2]{ht}{Hệ thống}
\newinst[2]{db}{Cơ sở dữ liệu}

\begin{call}{a}{1. Chọn xóa / ngừng kho}{gd}{Xác nhận}
\end{call}

\begin{call}{gd}{2. Gửi yêu cầu xóa}{ht}{}
  \begin{call}{ht}{3. Kiểm tra tồn / phiếu liên quan}{db}{Kết quả}
  \end{call}
  \begin{call}{ht}{4. Xóa hoặc ngừng Kho}{db}{Xác nhận}
  \end{call}
\end{call}

\begin{call}{ht}{5. Trả kết quả}{gd}{Thông báo}
\end{call}

\begin{call}{gd}{6. Cập nhật danh sách kho}{a}{}
\end{call}

\end{sequencediagram}
}
\caption{Sequence diagram UC20 -- Xóa kho}
\label{fig:sequence-uc20}
\end{figure}


Giải thích: Kiểm tra \texttt{TonKho}/phiếu trước khi xóa cứng hoặc ngừng.

%---------- UC21 ----------
\subsection{UC21 -- Tìm kiếm kho}

\textbf{Bối cảnh}: Tra cứu kho theo mã, tên, địa chỉ.

\begin{figure}[H]
\centering
\resizebox{\textwidth}{!}{%
\begin{sequencediagram}
\newinst{a}{Quản lý kho}
\newinst[2]{gd}{Giao diện kho}
\newinst[2]{ht}{Hệ thống}
\newinst[2]{db}{Cơ sở dữ liệu}

\begin{call}{a}{1. Nhập từ khóa tìm kho}{gd}{}
\end{call}

\begin{call}{gd}{2. Gửi điều kiện tìm}{ht}{}
  \begin{call}{ht}{3. Truy vấn Kho}{db}{Danh sách}
  \end{call}
\end{call}

\begin{call}{ht}{4. Trả kết quả}{gd}{}
\end{call}

\begin{call}{gd}{5. Hiển thị kho khớp}{a}{}
\end{call}

\end{sequencediagram}
}
\caption{Sequence diagram UC21 -- Tìm kiếm kho}
\label{fig:sequence-uc21}
\end{figure}


Giải thích: Query chỉ đọc \texttt{Kho}.

%---------- UC22 ----------
\subsection{UC22 -- Xem danh sách kho}

\textbf{Bối cảnh}: Xem toàn bộ danh mục kho.

\begin{figure}[H]
\centering
\resizebox{\textwidth}{!}{%
\begin{sequencediagram}
\newinst{a}{Quản lý kho}
\newinst[2]{gd}{Giao diện kho}
\newinst[2]{ht}{Hệ thống}
\newinst[2]{db}{Cơ sở dữ liệu}

\begin{call}{a}{1. Mở danh mục kho}{gd}{}
\end{call}

\begin{call}{gd}{2. Yêu cầu danh sách}{ht}{}
  \begin{call}{ht}{3. Lấy danh sách Kho}{db}{Dữ liệu}
  \end{call}
\end{call}

\begin{call}{ht}{4. Trả danh sách}{gd}{}
\end{call}

\begin{call}{gd}{5. Hiển thị bảng kho}{a}{}
\end{call}

\end{sequencediagram}
}
\caption{Sequence diagram UC22 -- Xem danh sách kho}
\label{fig:sequence-uc22}
\end{figure}


Giải thích: \texttt{layDanhSach} kho; điểm vào CRUD kho.

%---------- UC23 ----------
\subsection{UC23 -- Lập phiếu nhập kho}

\textbf{Bối cảnh}: NV lập phiếu (NCC, kho, chi tiết); gửi duyệt; QL duyệt tăng tồn.

\input{diagrams/tex/5.3_sequence_uc09}

Giải thích: Lưu phiếu + chi tiết (Chờ duyệt). Nhánh duyệt (có thể tách bước QL) gọi \texttt{tangTon}.

%---------- UC24 ----------
\subsection{UC24 -- Sửa phiếu nhập}

\textbf{Bối cảnh}: Sửa phiếu khi còn trạng thái cho phép (chưa duyệt).

\begin{figure}[H]
\centering
\resizebox{\textwidth}{!}{%
\begin{sequencediagram}
\newinst{a}{Nhân viên kho}
\newinst[2]{gd}{Giao diện phiếu nhập}
\newinst[2]{ht}{Hệ thống}
\newinst[2]{db}{Cơ sở dữ liệu}

\begin{call}{a}{1. Chọn phiếu (Chờ duyệt) để sửa}{gd}{Hiển thị form}
\end{call}

\begin{call}{a}{2. Sửa thông tin / chi tiết hàng}{gd}{}
\end{call}

\begin{call}{gd}{3. Gửi cập nhật phiếu}{ht}{}
  \begin{call}{ht}{4. Kiểm tra trạng thái cho phép sửa}{db}{Kết quả}
  \end{call}
  \begin{call}{ht}{5. Cập nhật PhieuNhap + ChiTiet}{db}{Xác nhận}
  \end{call}
\end{call}

\begin{call}{ht}{6. Trả kết quả}{gd}{Thành công}
\end{call}

\begin{call}{gd}{7. Hiển thị phiếu đã sửa}{a}{}
\end{call}

\end{sequencediagram}
}
\caption{Sequence diagram UC24 -- Sửa phiếu nhập}
\label{fig:sequence-uc24}
\end{figure}


Giải thích: Kiểm tra trạng thái $\rightarrow$ cập nhật header/chi tiết; \textbf{không} đổi tồn.

%---------- UC25 ----------
\subsection{UC25 -- Xóa phiếu nhập}

\textbf{Bối cảnh}: Xóa/hủy phiếu nhập chưa duyệt.

\input{diagrams/tex/5.seq_uc25}

Giải thích: Chỉ khi chưa duyệt; xóa master--detail, không đụng tồn.

%---------- UC26 ----------
\subsection{UC26 -- Tìm kiếm phiếu nhập}

\textbf{Bối cảnh}: Tìm theo mã phiếu, NCC, kho, khoảng ngày, trạng thái.

\begin{figure}[H]
\centering
\resizebox{\textwidth}{!}{%
\begin{sequencediagram}
\newinst{a}{Nhân viên / QL kho}
\newinst[2]{gd}{Giao diện phiếu nhập}
\newinst[2]{ht}{Hệ thống}
\newinst[2]{db}{Cơ sở dữ liệu}

\begin{call}{a}{1. Nhập điều kiện tìm phiếu}{gd}{}
\end{call}

\begin{call}{gd}{2. Gửi điều kiện tìm}{ht}{}
  \begin{call}{ht}{3. Truy vấn PhieuNhapKho}{db}{Danh sách}
  \end{call}
\end{call}

\begin{call}{ht}{4. Trả kết quả}{gd}{}
\end{call}

\begin{call}{gd}{5. Hiển thị phiếu khớp}{a}{}
\end{call}

\end{sequencediagram}
}
\caption{Sequence diagram UC26 -- Tìm kiếm phiếu nhập}
\label{fig:sequence-uc26}
\end{figure}


Giải thích: Query \texttt{PhieuNhapKho} theo bộ lọc.

%---------- UC27 ----------
\subsection{UC27 -- Xem danh sách phiếu nhập}

\textbf{Bối cảnh}: Xem danh sách phiếu nhập; chọn để sửa/duyệt/xem chi tiết.

\begin{figure}[H]
\centering
\resizebox{\textwidth}{!}{%
\begin{sequencediagram}
\newinst{a}{Nhân viên / QL kho}
\newinst[2]{gd}{Giao diện phiếu nhập}
\newinst[2]{ht}{Hệ thống}
\newinst[2]{db}{Cơ sở dữ liệu}

\begin{call}{a}{1. Mở danh sách phiếu nhập}{gd}{}
\end{call}

\begin{call}{gd}{2. Yêu cầu danh sách}{ht}{}
  \begin{call}{ht}{3. Lấy danh sách PhieuNhapKho}{db}{Dữ liệu}
  \end{call}
\end{call}

\begin{call}{ht}{4. Trả danh sách}{gd}{}
\end{call}

\begin{call}{gd}{5. Hiển thị bảng phiếu nhập}{a}{}
\end{call}

\end{sequencediagram}
}
\caption{Sequence diagram UC27 -- Xem danh sách phiếu nhập}
\label{fig:sequence-uc27}
\end{figure}


Giải thích: Tải danh sách phiếu; điểm vào UC23--UC26 và duyệt.

%---------- UC28 ----------
\subsection{UC28 -- Lập phiếu xuất kho}

\textbf{Bối cảnh}: Lập phiếu xuất; kiểm tra tồn khi thêm dòng; duyệt giảm tồn.

\input{diagrams/tex/5.4_sequence_uc10}

Giải thích: \texttt{kiemTraDu} trước khi nhận dòng; duyệt $\rightarrow$ \texttt{giamTon}.

%---------- UC29 ----------
\subsection{UC29 -- Sửa phiếu xuất}

\textbf{Bối cảnh}: Sửa phiếu xuất chưa duyệt; đổi SL có thể kiểm tra tồn lại.

\begin{figure}[H]
\centering
\resizebox{\textwidth}{!}{%
\begin{sequencediagram}
\newinst{a}{Nhân viên kho}
\newinst[2]{gd}{Giao diện phiếu xuất}
\newinst[2]{ht}{Hệ thống}
\newinst[2]{db}{Cơ sở dữ liệu}

\begin{call}{a}{1. Chọn phiếu (Chờ duyệt) để sửa}{gd}{Hiển thị form}
\end{call}

\begin{call}{a}{2. Sửa thông tin / chi tiết hàng}{gd}{}
\end{call}

\begin{call}{gd}{3. Gửi cập nhật phiếu}{ht}{}
  \begin{call}{ht}{4. Kiểm tra trạng thái + tồn (nếu đổi SL)}{db}{Kết quả}
  \end{call}
  \begin{call}{ht}{5. Cập nhật PhieuXuat + ChiTiet}{db}{Xác nhận}
  \end{call}
\end{call}

\begin{call}{ht}{6. Trả kết quả}{gd}{Thành công}
\end{call}

\begin{call}{gd}{7. Hiển thị phiếu đã sửa}{a}{}
\end{call}

\end{sequencediagram}
}
\caption{Sequence diagram UC29 -- Sửa phiếu xuất}
\label{fig:sequence-uc29}
\end{figure}


Giải thích: Tương tự UC24; thêm kiểm tra tồn nếu thay đổi số lượng xuất.

%---------- UC30 ----------
\subsection{UC30 -- Xóa phiếu xuất}

\textbf{Bối cảnh}: Xóa/hủy phiếu xuất chưa duyệt.

\begin{figure}[H]
\centering
\resizebox{\textwidth}{!}{%
\begin{sequencediagram}
\newinst{a}{Nhân viên kho}
\newinst[2]{gd}{Giao diện phiếu xuất}
\newinst[2]{ht}{Hệ thống}
\newinst[2]{db}{Cơ sở dữ liệu}

\begin{call}{a}{1. Chọn xóa / hủy phiếu xuất}{gd}{Xác nhận}
\end{call}

\begin{call}{gd}{2. Gửi yêu cầu xóa}{ht}{}
  \begin{call}{ht}{3. Kiểm tra chưa duyệt}{db}{Kết quả}
  \end{call}
  \begin{call}{ht}{4. Xóa/hủy PhieuXuat (+ chi tiết)}{db}{Xác nhận}
  \end{call}
\end{call}

\begin{call}{ht}{5. Trả kết quả}{gd}{Thông báo}
\end{call}

\begin{call}{gd}{6. Cập nhật danh sách phiếu}{a}{}
\end{call}

\end{sequencediagram}
}
\caption{Sequence diagram UC30 -- Xóa phiếu xuất}
\label{fig:sequence-uc30}
\end{figure}


Giải thích: Chỉ khi chưa duyệt; không hoàn tồn (vì chưa trừ).

%---------- UC31 ----------
\subsection{UC31 -- Tìm kiếm phiếu xuất}

\textbf{Bối cảnh}: Tìm phiếu xuất theo mã, kho, ngày, trạng thái.

\begin{figure}[H]
\centering
\resizebox{\textwidth}{!}{%
\begin{sequencediagram}
\newinst{a}{Nhân viên / QL kho}
\newinst[2]{gd}{Giao diện phiếu xuất}
\newinst[2]{ht}{Hệ thống}
\newinst[2]{db}{Cơ sở dữ liệu}

\begin{call}{a}{1. Nhập điều kiện tìm phiếu}{gd}{}
\end{call}

\begin{call}{gd}{2. Gửi điều kiện tìm}{ht}{}
  \begin{call}{ht}{3. Truy vấn PhieuXuatKho}{db}{Danh sách}
  \end{call}
\end{call}

\begin{call}{ht}{4. Trả kết quả}{gd}{}
\end{call}

\begin{call}{gd}{5. Hiển thị phiếu khớp}{a}{}
\end{call}

\end{sequencediagram}
}
\caption{Sequence diagram UC31 -- Tìm kiếm phiếu xuất}
\label{fig:sequence-uc31}
\end{figure}


Giải thích: Query \texttt{PhieuXuatKho}.

%---------- UC32 ----------
\subsection{UC32 -- Xem danh sách phiếu xuất}

\textbf{Bối cảnh}: Xem danh sách phiếu xuất.

\begin{figure}[H]
\centering
\resizebox{\textwidth}{!}{%
\begin{sequencediagram}
\newinst{a}{Nhân viên / QL kho}
\newinst[2]{gd}{Giao diện phiếu xuất}
\newinst[2]{ht}{Hệ thống}
\newinst[2]{db}{Cơ sở dữ liệu}

\begin{call}{a}{1. Mở danh sách phiếu xuất}{gd}{}
\end{call}

\begin{call}{gd}{2. Yêu cầu danh sách}{ht}{}
  \begin{call}{ht}{3. Lấy danh sách PhieuXuatKho}{db}{Dữ liệu}
  \end{call}
\end{call}

\begin{call}{ht}{4. Trả danh sách}{gd}{}
\end{call}

\begin{call}{gd}{5. Hiển thị bảng phiếu xuất}{a}{}
\end{call}

\end{sequencediagram}
}
\caption{Sequence diagram UC32 -- Xem danh sách phiếu xuất}
\label{fig:sequence-uc32}
\end{figure}


Giải thích: Tải danh sách; điểm vào CRUD/duyệt phiếu xuất.

%---------- UC33 ----------
\subsection{UC33 -- Lập phiếu kiểm kê}

\textbf{Bối cảnh}: NV tạo phiếu theo kho, nhập SL thực tế; QL duyệt gán lại tồn.

\begin{figure}[H]
\centering
\resizebox{\textwidth}{!}{%
\begin{sequencediagram}
\newinst{nv}{Nhân viên kho}
\newinst[1.5]{ql}{Quản lý kho}
\newinst[1.5]{gd}{Giao diện kiểm kê}
\newinst[1.5]{ht}{Hệ thống}
\newinst[1.5]{db}{CSDL}

\begin{call}{nv}{1. Tạo phiếu kiểm kê (chọn kho)}{gd}{}
\end{call}

\begin{call}{gd}{2. Yêu cầu tạo phiếu}{ht}{}
  \begin{call}{ht}{3. Nạp SL sổ từ TonKho}{db}{Danh sách dòng}
  \end{call}
  \begin{call}{ht}{4. Lưu PhieuKiemKe + chi tiết}{db}{Xác nhận}
  \end{call}
\end{call}

\begin{call}{nv}{5. Nhập số lượng thực tế}{gd}{}
\end{call}

\begin{call}{gd}{6. Cập nhật chi tiết, tính chênh lệch}{ht}{}
  \begin{call}{ht}{7. Lưu ChiTietKiemKe}{db}{}
  \end{call}
\end{call}

\begin{call}{nv}{8. Gửi duyệt}{gd}{}
\end{call}

\begin{call}{ql}{9. Duyệt điều chỉnh tồn}{gd}{}
\end{call}

\begin{call}{gd}{10. Xác nhận duyệt}{ht}{}
  \begin{call}{ht}{11. ganTheoKiemKe (cập nhật TonKho)}{db}{Xác nhận}
  \end{call}
\end{call}

\begin{call}{ht}{12. Thông báo hoàn tất}{gd}{}
\end{call}

\end{sequencediagram}
}
\caption{Sequence diagram UC33 -- Lập phiếu kiểm kê}
\label{fig:sequence-uc33}
\end{figure}


Giải thích: Nạp SL sổ từ \texttt{TonKho} $\rightarrow$ nhập thực tế $\rightarrow$ chênh lệch $\rightarrow$ duyệt mới \texttt{ganTheoKiemKe}.

%---------- UC34 ----------
\subsection{UC34 -- Xem báo cáo nhập xuất tồn}

\textbf{Bối cảnh}: QL tổng hợp NXT theo kỳ/kho/hàng; có thể xuất Excel.

\begin{figure}[H]
\centering
\resizebox{\textwidth}{!}{%
\begin{sequencediagram}
\newinst{ql}{Quản lý kho}
\newinst[2]{gd}{Giao diện báo cáo}
\newinst[2]{ht}{Hệ thống}
\newinst[2]{db}{Cơ sở dữ liệu}

\begin{call}{ql}{1. Mở báo cáo NXT}{gd}{Form điều kiện}
\end{call}

\begin{call}{ql}{2. Chọn kỳ, kho, hàng (tùy chọn)}{gd}{}
\end{call}

\begin{call}{gd}{3. Gửi điều kiện báo cáo}{ht}{}
  \begin{call}{ht}{4. Tổng hợp phiếu nhập đã duyệt}{db}{Dữ liệu nhập}
  \end{call}
  \begin{call}{ht}{5. Tổng hợp phiếu xuất đã duyệt}{db}{Dữ liệu xuất}
  \end{call}
  \begin{call}{ht}{6. Lấy tồn hiện tại (TonKho)}{db}{Dữ liệu tồn}
  \end{call}
\end{call}

\begin{call}{ht}{7. Trả bảng NXT}{gd}{}
\end{call}

\begin{call}{gd}{8. Hiển thị / xuất Excel}{ql}{}
\end{call}

\end{sequencediagram}
}
\caption{Sequence diagram UC34 -- Xem báo cáo nhập xuất tồn}
\label{fig:sequence-uc34}
\end{figure}


Giải thích: Chỉ \textbf{đọc/tổng hợp} phiếu đã duyệt và tồn; không tạo giao dịch mới.

\section{Phân tích hành vi đối tượng phiếu (biểu đồ trạng thái)}

Sequence mô tả tương tác \textbf{giữa} các đối tượng; state machine mô tả hành vi \textbf{của chính} đối tượng phiếu qua các trạng thái.

\subsection{Phiếu nhập kho}

Trạng thái điển hình: Mới tạo $\rightarrow$ Chờ duyệt $\rightarrow$ Đã duyệt / Từ chối; có thể Hủy khi được phép. Chỉ ``Đã duyệt'' mới làm thay đổi tồn kho.

\input{diagrams/tex/5.5_state_phieu_nhap}

Giải thích: Sự kiện chuyển: gửi duyệt, duyệt, từ chối, hủy. ``Đã duyệt'' là điều kiện tăng tồn.

\subsection{Phiếu xuất kho}

Cùng mô hình trạng thái với phiếu nhập.

\input{diagrams/tex/5.6_state_phieu_xuat}

Giải thích: Đối xứng phiếu nhập; khi ``Đã duyệt'' thì \textbf{giảm} tồn.

\section{Kết luận chương}

Chương 5 hoàn thành \textbf{phân tích hành vi}:
\begin{itemize}[leftmargin=1.5cm]
    \item \textbf{34 biểu đồ trình tự} tương ứng \textbf{34 use case} (UC01--UC34);
    \item State bổ sung vòng đời phiếu nhập/xuất.
\end{itemize}

Tổng kết ba giai đoạn:
\begin{enumerate}[leftmargin=1.5cm]
    \item \textbf{Chức năng} (Ch.3): Use Case;
    \item \textbf{Cấu trúc} (Ch.4): Class;
    \item \textbf{Hành vi} (Ch.5): Sequence (+ State).
\end{enumerate}

Tiếp theo: thiết kế chi tiết (Ch.6) và CSDL/giao diện (Ch.7).
